% =============================================================
%  Kapitel 4 – Dynamisches Verhalten (Global Software Control + Boundary Conditions)
% =============================================================
\section{Dynamisches Verhalten des Systems}
\label{sec:dynamik}

Dieser Abschnitt beschreibt das Laufzeitverhalten anhand von Sequenzdiagrammen.
Er umfasst sowohl den \emph{Global Software Control} (regulärer Betrieb mit
Emergency Stop) als auch die \emph{Boundary Conditions} (Systemstart, Shutdown,
Fehlerbehandlung).

% ------------------------------------------------------------
\subsection{Global Software Control: Auslösung durch Safety-Zone}
\label{sec:seq_safetyzone}

\begin{figure}[H]
  \centering
  \includegraphics[width=0.95\textwidth]{04_Sequence_SafetyZone}
  \caption{Sequenzdiagramm: Auslösung des Emergency Stop durch Safety-Zone (Rot)}
  \label{fig:seq_safetyzone}
\end{figure}

Wird eine Distanz $\leq 15\,\mathrm{cm}$ erkannt (rote Sicherheitszone), läuft
folgender Ablauf ab (Abbildung~\ref{fig:seq_safetyzone}):

\begin{enumerate}
  \item Der \texttt{SensorFusionAdapter} auf dem Jetson meldet ein Objekt in der
        Safety-Zone.
  \item Der \texttt{TTCObserver} erzeugt ein \texttt{ZoneUpdate(Safety)}-Ereignis.
  \item Über \ac{LAN} $\to$ \texttt{MasterLanLink} $\to$ \texttt{MasterCanGateway}
        wird ein \ac{CAN}-Paket generiert.
  \item Das \ac{CAN}-Paket wird vom \texttt{EnvironmentControl}-Interface auf dem Slave
        empfangen und an \texttt{EmergencyStop} weitergegeben.
  \item Der Emergency Stop wechselt in den Zustand \texttt{EmergencyStop},
        sendet einen \textbf{Stop-Befehl} an die \texttt{DriveControlUnit}
        und löst einen \texttt{Logger.log(ERROR, "Safety zone triggered")}-Aufruf aus.
  \item Die Status-\ac{LED} wechselt auf Rot.
\end{enumerate}

\textbf{Wichtig:} Das Fahrzeug bleibt blockiert, bis eine manuelle Freigabe
die Guard-Bedingung erfüllt (vgl.\ Abschnitt~\ref{sec:seq_release}).

% ------------------------------------------------------------
\subsection{Boundary Condition: Systemstart (Startup)}
\label{sec:seq_startup}

\begin{figure}[H]
  \centering
  \includegraphics[width=0.95\textwidth]{05_Sequence_Startup}
  \caption{Sequenzdiagramm: Startup-Verhalten}
  \label{fig:seq_startup}
\end{figure}

Beim Systemstart (Abbildung~\ref{fig:seq_startup}) werden alle Komponenten initialisiert.
Der Emergency Stop beginnt automatisch im Zustand \texttt{SoftStop}
(\emph{Safe-Default}): Das Fahrzeug ist erst nach einer expliziten
manuellen Freigabe fahrbereit.

\textbf{Ablauf:}
\begin{enumerate}
  \item Alle Hardware-Einheiten starten ihre Initialisierungsroutine.
  \item \texttt{EmergencyStop} wechselt von \texttt{Initialization}
        nach \texttt{SoftStop}.
  \item \texttt{Logger} registriert \texttt{INFO: System startup complete}.
  \item Fahrzeug wartet auf Freigabe-Signal (Taster/Button).
\end{enumerate}

% ------------------------------------------------------------
\subsection{Boundary Condition: Freigabe nach Emergency Stop (Release)}
\label{sec:seq_release}

\begin{figure}[H]
  \centering
  \includegraphics[width=0.95\textwidth]{05_Sequence_Release}
  \caption{Sequenzdiagramm: Manuelle Freigabe nach Emergency Stop}
  \label{fig:seq_release}
\end{figure}

Nach einem Emergency Stop ist eine manuelle Freigabe erforderlich (Abbildung~\ref{fig:seq_release}).
Die Guard-Bedingung muss vollständig erfüllt sein:

\begin{itemize}
  \item Safety-Zone ist \textbf{nicht} mehr aktiv (kein Objekt in Rot).
  \item Kein Fehler ist aktiv (\texttt{errorFlag == false}).
\end{itemize}

Erst dann erfolgt:
\begin{enumerate}
  \item Zustandswechsel \texttt{EmergencyStop} $\to$ \texttt{Normal}.
  \item Aktivierung der \texttt{DriveControlUnit} (Freigabe der Motorsteuerung).
  \item \texttt{Logger.log(INFO, "System released, resuming normal operation")}.
\end{enumerate}

% ------------------------------------------------------------
\subsection{Boundary Condition: Freigabe abgelehnt (Release Rejected)}
\label{sec:seq_releaserejected}

\begin{figure}[H]
  \centering
  \includegraphics[width=0.95\textwidth]{05_Sequence_ReleaseRejected}
  \caption{Sequenzdiagramm: Freigabe abgelehnt (Guard nicht erfüllt)}
  \label{fig:seq_releaserejected}
\end{figure}

Ist die Safety-Zone noch aktiv oder liegt ein Fehler vor,
wird der Freigabe-Request \textbf{abgelehnt} (Abbildung~\ref{fig:seq_releaserejected}).
Der Zustand bleibt unverändert (\texttt{EmergencyStop}).
\texttt{Logger} protokolliert \texttt{WARNING: Release rejected – condition not met}.

% ------------------------------------------------------------
\subsection{Boundary Condition: Fehlerbehandlung (Error)}
\label{sec:seq_error}

\begin{figure}[H]
  \centering
  \includegraphics[width=0.95\textwidth]{05_Sequence_Error}
  \caption{Sequenzdiagramm: Fehlerbehandlung (fehlende Sensordaten)}
  \label{fig:seq_error}
\end{figure}

Fehlen gültige Sensordaten (z.\,B.\ LiDAR-Ausfall, \ac{CAN}-Timeout),
greift das \textbf{Fail-Safe-Prinzip} (Abbildung~\ref{fig:seq_error}):

\begin{enumerate}
  \item \texttt{EnvironmentControl} detektiert fehlende/ungültige \ac{CAN}-Daten.
  \item \texttt{EmergencyStop} erhält \texttt{ErrorDetected}-Ereignis.
  \item Zustandswechsel nach \texttt{SoftStop} (falls noch nicht gestoppt).
  \item \texttt{Logger.log(ERROR, "Sensor data missing – soft stop engaged")}.
\end{enumerate}

\textbf{Regel:} Keine gültigen Daten $\Rightarrow$ Stopp + Logging (keine Ausnahmen).

% ------------------------------------------------------------
\subsection{Boundary Condition: Shutdown}
\label{sec:seq_shutdown}

\begin{figure}[H]
  \centering
  \includegraphics[width=0.95\textwidth]{05_Sequence_Shutdown}
  \caption{Sequenzdiagramm: Kontrollierter Shutdown}
  \label{fig:seq_shutdown}
\end{figure}

Bei einem \texttt{ShutdownRequest} (Abbildung~\ref{fig:seq_shutdown}):
\begin{enumerate}
  \item \texttt{EmergencyStop} empfängt das Shutdown-Ereignis.
  \item Zustandswechsel nach \texttt{SoftStop}.
  \item Stop-Befehl an \texttt{DriveControlUnit}.
  \item \texttt{Logger.log(INFO, "Shutdown initiated")}.
  \item Alle Subsysteme fahren geordnet herunter.
\end{enumerate}
